\documentclass[]{article}

%\usepackage{epcc}
\usepackage{jmlr2e}
%\usepackage[hyphens]{url}
%\PassOptionsToPackage{hyphens}{url}
%\usepackage[colorlinks=false, hidelinks]{hyperref}
\usepackage{amsmath}
\usepackage[toc,page]{appendix}
\usepackage[table]{xcolor}
\usepackage[marginparsep=30pt]{geometry}
\usepackage{stmaryrd}
\usepackage{algorithm}
\usepackage{algorithmic}
\usepackage{tikz}
\usepackage{pgfplots}
\usepackage{tabu}
\usepackage{longtable}
\usepackage{tabularx}
\usepackage{listings}
\usepackage{fancyref}
\usepackage{relsize}
\usepackage{float}
\usepackage{graphicx}
\usepackage{subcaption}
\usepackage{diagbox}
\usepackage{multirow}
\usepackage{slashbox}
\usepackage{graphics}
\usepackage{booktabs}
\usepackage{natbib}
\usepackage{csquotes}

\usetikzlibrary{%
    arrows,
    arrows.meta,
    decorations,
    backgrounds,
    positioning,
    fit,
    petri,
    shadows,
    datavisualization.formats.functions,
    calc,
    shapes,
    shapes.multipart,
    matrix,
    plotmarks
}

\usepgfplotslibrary{fillbetween, statistics, dateplot}

\pgfplotsset{%
  compat=1.3,
  every non boxed x axis/.style={%
  enlarge x limits=false,
  x axis line style={}%-stealth},
  },
  every boxed x axis/.style={},
  every non boxed y axis/.style={%
  enlarge y limits=false,
  y axis line style={}%-stealth},
  },
  every boxed y axis/.style={},
}

\renewcommand{\labelenumii}{\theenumii}
\renewcommand{\theenumii}{\theenumi.\arabic{enumii}.}

\bibliographystyle{plainnat}
\bibpunct{(}{)}{;}{a}{,}{,}

\newenvironment{declaration}
{\centerline{\large\bf Declaration}\vspace{0.7ex}%
  \bgroup\leftskip 20pt\rightskip 20pt\noindent\ignorespaces}%
{\par\egroup\vskip 0.25ex}


\def\title{Deep Learning on SpiNNaker}
\def\author{Jonas Fassbender \\ \textit{jonas@fassbender.dev}}
\date{}

\ShortHeadings{Jonas Fassbender}{\title}

\begin{document}

% titlepage {{{
\begin{titlepage}

\begin{flushleft}
	\vspace*{-1cm}
	\includegraphics[scale=0.15]{logos/logo_black.pdf}\\
	\vspace*{1cm}
\end{flushleft}
\begin{flushright}
	\vspace*{-3cm}
	\includegraphics[scale=0.2]{logos/crest_bw.pdf}\\
	\vspace*{1cm}
\end{flushright}


%\EPCCheaderLogos
\null

\begin{center}
\begin{Huge}
\textbf{\title}
\end{Huge}
~\\
~\\
~\\
\textit{\Large {\LARGE M}ASTER {\LARGE T}HESIS}
~\\
~\\
~\\
\begin{Large}
\begin{tabu} to \textwidth {Xr}
Jonas Fassbender
&\href{mailto:jonas@fassbender.dev}{jonas@fassbender.dev}
\end{tabu}
\end{Large}
~\\
~\\
~\\
\begin{large}
In the course of studies

\textit{{\Large H}IGH {\Large P}ERFORMANCE {\Large C}OMPUTING WITH {\Large D}ATA {\Large S}CIENCE}
~\\
~\\
~\\
For the degree of

\textit{{\Large M}ASTER OF {\Large S}CIENCE}
~\\
~\\
~\\
The University of Edinburgh
~\\
~\\
~\\
\begin{tabular}{rl}
  First supervisor: &Caoimhín Laoide-Kemp \\
                    &EPCC \\
  &\\
  Second supervisor: &Dr Kevin Stratford \\
                     &EPCC \\
  &\\
  Third supervisor: &Dr Alan Stokes \\
                    &University of Manchester \\
\end{tabular}
~\\
~\\
~\\
Edinburgh, August 2020
\end{large}
\end{center}
\end{titlepage}
% }}}

\pagenumbering{roman}

% here thanks

% declaration {{{
\hspace{0pt}
\vfill

\begin{declaration}
I declare that this dissertation was composed by myself, that the work
contained herein is my own except where explicitly stated otherwise in
the text, and that this work has not been submitted for any other
degree or professional qualification except as specified.
~\\
~\\
~\\
\begin{tabu}{Xc}
  &Jonas Fassbender \\
  &August 2020
\end{tabu}
\end{declaration}

\vfill
\hspace{0pt}
% }}}

\newpage

% abstract {{{
\hspace{0pt}
\vfill

\begin{abstract}
\end{abstract}

\vfill
\hspace{0pt}
% }}}

\newpage

\tableofcontents

\newpage

% listoffigures

% listoftables

\pagenumbering{arabic}

% fancy page (maybe above arabic -- depends on page numbers)

\section{Introduction}

Deep learning is revolutionizing the world.
It has become part of our daily lives as consumers, powering major
software products---from recommendation systems over translation tools
to web search \citep{lecun_et_al_2015}.
Major breakthroughs in fields like computer vision or natural language
processing were achieved through the use of deep learning
\citep{krizhevsky_et_al_2012, hinton_et_al_2012}.
It has emerged as a driving force behind discoveries in numerous
domains like particle physics, drug discovery, genomics or gaming
\citep{ciodaro_et_al_2012, ma_et_al_2015, leung_et_al_2014,
  silver_et_al_2016}.

Deep learning has become so ubiquitous that we are changing the
way we build modern hardware to account for its computational demands.
From the way edge devices like mobile phones or embedded systems are
designed over modern CPUs to specialized hardware for deep learning
models \citep{TODO}.
Whole supercomputers are build solely for deep learning \citep{TODO}.

Another field besides deep learning has high computational demands for
specific algorithms: neuroscience.

\section{Background}

\subsection{An Introduction to Deep Learning}

% history
% neural network as a statistical method, not biologically correct

% non-linear classifiers

% concepts: layers, activation, optimizer, forward- and backward-pass
%           hidden units, ... (chapter 6 goodfellow)

% CNN's (cnn, relu, pooling, ...)

% (Residual stuff)

\subsection{Computer Vision: ImageNet and the ILSVRC}

% short section describing imagenet and ilsvrc and their importance
% for image recognition (computer vision) tasks

% ResNet50

\subsection{Benchmarking Deep Learning Systems: The MLPerf Benchmark}

% also rather short

% what is the point

% some benchmarks

% if references are too small: add all the papers from the report

\subsection{SpiNNaker as a Neuromorphic Computer Architecture}

% spinnaker architecture

% compare to other state of the art machine learning hardware

% why didn't we just implement a backend (e.g. tensorflow?)
% because architecture is so different (MACs vs. massive parallelism
% with small messages)

\subsection{Related Work}

% that neural thing that translates trained models to spiking nns

% find more about DL on SpiNNaker


\section{Deep Learning on SpiNNaker: SpiDNN}

\section{Benchmark}

\section{Discussion}

\section{Conclusion}

\bibliography{library.bib}
\addcontentsline{toc}{section}{References}

\end{document}
